\section{Case Study}
\label{sec:casestudy}

\TODO{Pierre: The current section explains the vulnerability well, but it does not yet function as a case study of ToFU. The key purpose of this section should be to walk through the complete discovery process. Explain what metadata ToFU receives for this tool, which sanitizer is selected and why, what input or seed is generated, how fuzzing transforms that input into the triggering argument, what runtime operation the introspection layer observes, and why the semantic oracle identifies that operation as a violation. The reader should be able to follow the bug from tool metadata all the way to the final vulnerability.}

\TODO{Pierre: At the beginning of the case study, briefly explain why this particular bug is representative of ToFU's contribution. The interesting point is not simply that command-line argument injection exists, but that the argument is syntactically a valid string and is passed safely without a shell, while its semantic interpretation by the downstream command violates the intended tool contract. This is exactly the distinction the paper should emphasize.}

An example of a vulnerability found using \project{} that can be studied can be found in \texttt{tools-main/\allowbreak python/\allowbreak extract\_issues.py}. 

\TODO{Pierre: Before the code listing, describe the tool's documented interface. State what range\_spec is supposed to represent according to the tool metadata or description and, if available, include the relevant description supplied to ToFU. This establishes the intended semantic contract before showing how the implementation violates it.}

\begin{lstlisting}[
  basicstyle=\ttfamily\small,
  breaklines=true,
  columns=fullflexible,
  linewidth=\linewidth
]
def extract_issue_numbers(range_spec: str) -> str:
    ...
    git_range = range_spec if ".." in range_spec else f"{range_spec}..HEAD"
    ...
    result = subprocess.run(
        ["git", "log", git_range, "--pretty=%B"],
        capture_output=True,
        text=True,
        check=True,
    )
    ...

\end{lstlisting}

\TODO{Pierre: Immediately after the listing, connect this code to ToFU rather than moving directly to manual vulnerability analysis. Explain which sanitizer ToFU associates with range\_spec and why the tool metadata makes that sanitizer applicable. If ToFU specializes a generic CWE-88 sanitizer to this argument, show the resulting check or summarize it here. This is the most important missing step in the current case study.}

In this case, range\_spec is expected to be a string representing either a single tag/commit or an explicit Git range.

\TODO{Pierre: Explain how this expectation is known. Ideally tie it directly to the tool description, argument metadata, documentation, or generated sanitizer. This is important because ToFU's claim is that the security property follows from documented tool semantics rather than a manually written specification.}

By default, \texttt{subprocess.run} is invoked with a list of arguments and without a shell interpreter (\texttt{shell=False}); therefore, classic command chaining and shell injection (e.g., providing \texttt{range\_spec="; echo success"} or piping operators) are inherently prevented.

\TODO{Pierre: This distinction is important and should be connected explicitly to the motivation for semantic sanitization. A conventional command-injection check looking for shell metacharacters would consider this invocation substantially safer because shell interpretation is disabled. ToFU instead needs to reason about the semantics of the argument accepted by Git itself. Make that point explicit.}

However, passing unsanitized arguments directly into a command invocation opens the door to \textit{CLI Argument Injection} (CWE-88). Because the argument is passed as a standalone vector element to \texttt{git}, an input formatted as a command-line flag (such as \texttt{--output=/user/../user/foo}) causes the \texttt{git} binary to interpret the parameter as an operational option rather than a positional revision range. Under this flag, \texttt{git log} redirects its formatted output directly to the specified file path on disk.

\TODO{Pierre: Show the actual triggering input produced by ToFU, preferably exactly as recorded by the fuzzer after minimization. If the example shown here is the actual minimized input, state that explicitly. If it is only an illustrative input, replace it with or separately show the real input. A case study should distinguish the discovered proof of concept from an explanatory example.}

\TODO{Pierre: Explain how ToFU's runtime oracle detects this behavior. For example, if the introspection layer observes the file creation or the subprocess arguments, state exactly what event is intercepted and what rule causes ToFU to flag it. The current text explains why the behavior is vulnerable, but not how ToFU automatically recognizes the violation.}

\TODO{Pierre: Also explain whether ToFU recognizes the dangerous command-line option before execution or detects its effect afterward through runtime introspection. This distinction matters for understanding the architecture and should match the Design section.}

Furthermore, the ternary check on line 3 (\texttt{range\_spec if ".." in range\_spec else f"\{range\_spec\}..HEAD"}) is subverted by directory traversal sequences: providing path traversal dots (\texttt{..}) in the file path satisfies the substring condition, thereby bypassing the automatic appending of \texttt{..HEAD} and passing the adversarial \texttt{--output} argument verbatim to the subprocess on line 6.

\TODO{Pierre: There appears to be a LaTeX or code-formatting issue in the inline expression in this sentence. Check the nested texttt commands and braces carefully. I am not changing it here because the existing text should remain untouched, but this line should be fixed before submission.}

\TODO{Pierre: Clarify that the two dots are useful here because the implementation checks only whether the string contains two consecutive dots. The security problem is not conventional path traversal by itself. The dots also manipulate the program's branch condition so that the injected Git option is passed unchanged. That interaction is an interesting semantic aspect of the bug and is worth emphasizing.}

\TODO{Pierre: If ToFU generated or mutated the two-dot sequence automatically, explain that transition here. Showing how the fuzzer reaches the final option while satisfying the program's range check would make this a much stronger fuzzing case study.}

Although this vulnerability does not provide direct arbitrary code execution, it presents serious security and operational risks that can be exploited by an attacker in two distinct ways:

\noindent
\textbf{Denial of Service through Storage and Inode Exhaustion (CWE-400):} By repeatedly invoking the tool with distinct file names or deeply nested paths across the filesystem, an attacker can generate a large volume of output files. This continuous file generation can quickly saturate available disk space or exhaust filesystem inodes, ultimately causing denial of service for the MCP server and co-located services.

\noindent
\textbf{Arbitrary File Overwrite and Asset Integrity Corruption (CWE-732):} Depending on the filesystem permissions granted to the MCP server process, an attacker can overwrite existing critical files, such as configuration files, local database snapshots, or application assets. The overwritten files will contain the formatted \texttt{git log} commit history (\texttt{--pretty=\%B}) rather than their original contents. Consequently, any subsequent process or service attempting to read or parse these corrupted files will encounter unexpected structures and fail, inducing silent runtime errors, degraded application behavior, or complete subsystem crashes.


\TODO{Pierre: Recheck the CWE classification for the arbitrary-file-overwrite consequence. CWE-732 is Incorrect Permission Assignment for Critical Resource, which does not appear to describe the root cause here because the vulnerability does not assign incorrect file permissions. CWE-88 already captures the argument-injection root cause. If you want a secondary CWE for the resulting filesystem behavior, select one that actually corresponds to uncontrolled file path or file modification rather than permission assignment.}

\TODO{Pierre: Be careful with the claim about deeply nested paths. State only behavior that was actually validated. Git can write to a supplied output path, but this does not necessarily imply that it creates arbitrary missing parent directories. The denial-of-service argument does not need this stronger claim because repeatedly creating distinct files in writable existing directories is sufficient to explain storage or inode exhaustion.}

\TODO{Pierre: Close the case study with a short paragraph explaining why this bug illustrates ToFU's broader value. The key observation is that the parameter is type-correct, the subprocess invocation avoids shell interpretation, and ordinary malformed-input checking may therefore see nothing unusual. The vulnerability appears only when the string is interpreted according to Git's command-line semantics and compared against the intended meaning of range\_spec. Connect this explicitly back to the LLM-to-tool trust boundary and semantic sanitizer design.}

% TODO(Nazmul): Later, replay this vulnerability under a conventional Atheris
% or schema-guided fuzzing configuration without ToFU's semantic seeding and
% sanitizer guidance. Report whether the baseline discovers the argument
% injection under the same budget and compare time to discovery.

% TODO(Nazmul): Later, separately disable the semantic runtime oracle for this
% case study and determine whether the triggering input would otherwise appear
% as an ordinary successful tool execution. This would provide a particularly
% clear demonstration of why semantic runtime checking is necessary.

% TODO(Nazmul): If the original fuzzing logs retain the mutation history or
% corpus evolution, reconstruct the progression from the initial seed to the
% minimized --output triggering input. This could provide a useful concrete
% example of how ToFU's seeding and coverage-guided mutation cooperate.

